\documentclass[11pt]{article}
\usepackage{icmlko}
\begin{document}

\title{싼 자는 한쪽만 증명한다 --- 적합도를 안 사면 안전하고, 사면
알 수 없다}
\author{월드모델 실험실}
\date{노트 475 · 2026년 8월 3일}
\maketitle

\begin{abstract}
노트 474 가 조항을 하나 만들었다 --- \emph{용량 변경이 판에서 비기면 다른
자 최소 둘로 확인한다.} 옳지만 \textbf{후보마다 40분}이 든다. 그래서
\textbf{값싼 앞자}를 찾았다: 학습 적합도와 원핫 의존도 --- 둘 다 유보
라벨도 집 밖 짝도 안 쓴다.

등록한 문장은 \emph{``판이 비기는데 적합도가 크게 오르면 위험''} 이었다.

\begin{center}
\begin{tabular}{lrrl}
\hline
후보 & 적합도 차 & 원핫 의존 차 & 여섯 자 판정 \\
\hline
\texttt{it} $400$ & $\mathbf{+0.0434}$ & $+0.0016$ & \textbf{다섯에서 잃음} \\
lr $0.16$ & $\mathbf{+0.0351}$ & $+0.0009$ & 여섯 다 비김 \\
K $64$ & $+0.0001$ & $+0.0002$ & 여섯 다 비김 \\
깊이 $24$ & $+0.0007$ & $-0.0002$ & 비김 \\
\hline
\end{tabular}
\end{center}

\textbf{반증이다.} lr $0.16$ 이 \texttt{it} $400$ 의 $81\%$ 만큼 적합도를
사는데 \textbf{여섯 자에서 무해}하다. 적합도 상승만으로는 둘을 못 가른다.
원핫 의존도도 $1.8$배 차이라 마찬가지다.

\textbf{그런데 표가 한쪽으로는 깨끗하다.}
\end{abstract}

\begin{figure}[t]
\centering
\includegraphics[width=.86\textwidth]{figs/f1.pdf}
\caption{초록 칸(적합도를 안 사는 쪽)에 든 둘은 \emph{둘 다} 무해했고,
빨강 칸에 든 둘은 하나만 해롭다.}
\end{figure}

\section{반쪽짜리지만 쓸모가 있다}

\begin{quote}
\textbf{적합도를 안 사면 안전하다} --- K $64$($+0.0001$)와 깊이
$24$($+0.0007$)는 여섯 자에서 아무 일도 안 했다. \\
\textbf{적합도를 사면 알 수 없다} --- \texttt{it} $400$($+0.0434$)은
다섯에서 잃고 lr $0.16$($+0.0351$)은 안 잃는다.
\end{quote}

논리로 적으면 \emph{적합도 상승 $\approx 0$} 은 무해함의 \textbf{충분
조건처럼 보이고}(관측 $2/2$), 해로움의 \textbf{필요 조건}이다 --- 판을
속이려면 뭔가는 더 맞혀야 하니까.

\textbf{그래서 앞자로 쓸 수 있다. 다만 통과가 아니라 건너뛰기로.}

\begin{quote}
\textbf{조항} --- 용량 변경의 \textbf{학습 적합도 상승이 잡음 크기면
(여기서는 $+0.001$ 아래) 노트 474 의 비싼 확인을 건너뛴다.} 그보다 크면
\textbf{반드시} 비싼 확인을 한다 --- 싼 자는 그 안에서 아무것도 못
가른다.
\end{quote}

$2/2$ 는 적은 근거다. 그래서 \emph{건너뛰기}로만 쓰고, 건너뛴 변경은
장부에 ``싼 자로만 봤음''이라 적어 둔다. 반례가 하나 나오면 조항을
버린다.

\section{왜 lr 은 사고도 안 해로운가}

이것이 이 노트가 못 푼 것이다. \texttt{it} 과 lr 은 노트 $459 \sim 461$
에서 계속 같은 무리로 묶였는데(용량 가족 · 환율 $0.62 \sim 0.54$),
\textbf{여섯 자로 보면 다르다}. lr $0.16$ 은 적합도를 $+0.035$ 사고
\emph{아무 데서도 안 잃는다} --- 산 것이 어디로 갔는지 이 노트의 자로는
안 보인다.

가설 하나만 적어 둔다. lr 은 \emph{보폭}이라 같은 나무를 더 멀리 밀 뿐
구조를 안 바꾸고, \texttt{it} 은 \emph{나무 수}라 새 갈래를 계속 만든다.
새 갈래가 무리 사이를 파고들 자리다. \textbf{안 쟀다.}

\section*{남은 것}

\begin{itemize}
\item \textbf{lr 이 산 적합도는 어디 있나} --- 여섯 자 어디에도 안 나온다.
자가 모자란 것인지 진짜 아무 데도 안 가는 것인지.
\item \textbf{건너뛰기 조항의 반례 찾기} --- 적합도를 $+0.001$ 아래로
사면서 여섯 자 중 하나를 잃는 변경이 있나. 있으면 조항이 죽는다.
\item \textbf{축 추가에도 걸리나} --- 지금 조항은 용량 변경에만 걸었다.
\end{itemize}

\end{document}
